\section{Orbits of the multiplier group}\label{sec:orb}
\begin{proposition}\label{prop:orb}
Let $q=p^k$ and let $N(\rho)$ count triples $(a,b,c)\in S^3$ with $b-a$ a unit and $c-a=\rho(b-a)$ (for prime $m$ this is the count of Section~3).
Then $3k\mid N(\rho)$ for every $\rho$.
\end{proposition}
\begin{proof}
By the multiplier theorem, some translate of $S$ is fixed by multiplication by $p$ \cite{Ha,MR}. $N$ is translation invariant, so assume $pS=S$.
We have $p^{3k}-1=(q-1)m$. If $p^j\equiv1\pmod m$ with $0<j<3k$, then $m\mid p^d-1$ for $d=\gcd(j,3k)\le 2k$, which is impossible because $0<p^d-1<m$.
So $p$ has order $3k$ modulo $m$. The map $(a,b,c)\mapsto(pa,pb,pc)$ preserves the set of triples. A triple fixed by $p^j$, $0<j<3k$, would give
$(p^j-1)(b-a)=0$ with $b-a$ a unit, which is impossible. So the action is free and every orbit has $3k$ elements.
\end{proof}
This is elementary and probably known. The order statement and the freeness step are machine-checked in Lean~4 with Mathlib
(files \texttt{C9Order.lean}, \texttt{C9Orbit.lean}; standard axioms only). The multiplier theorem itself is not formalised.
So Conjecture~\ref{conj:var} says that $M=N/3k$ is Poisson-like. This holds for non-cube $q$ (Table~\ref{tab:pp}) but fails for cube $q$; the next two sections explain most of this failure.

\section{A resonance theorem for cube $q$}\label{sec:res}
Let $q=r^3$, where $r$ is a prime power. Let $F=\mathbb F_q$ and $K=\mathbb F_{q^3}=\mathbb F_{r^9}$. Write $\tau(u)=u^r$ and $\varphi=\tau^3$ (the $q$-Frobenius).
Let $L=\ker\mathrm{Tr}_{K/F}$, so that $S=\{i: g^i\in L\}$. Since $\tau$ commutes with $\varphi$, we have $\tau(L)=L$. Moreover $\varphi$ is $F$-linear on $L$ with $\varphi^3=1$.
Since $r^9=q^3\equiv1$ and $r^3=q\not\equiv1\pmod m$, $r$ has order $9$ modulo $m$. Put $\rho_1=r+r^{-1}\bmod m$.
\begin{theorem}\label{thm:res}
Let $x\in P(L)$ be a point of the Singer line.
\begin{itemize}
\item If $\varphi x\ne x$ (projectively), then there are exactly $r$ points $y\in P(L)\setminus\{x\}$ with $z/x=(y/x)^{\rho_1}$ for some $z\in P(L)$.
\item If $\varphi x=x$, there are none.
\end{itemize}
The number $e$ of $\varphi$-fixed points of $P(L)$ is $0$, $1$ or $2$ according as $r\equiv2$, $0$ or $1\pmod3$.
Consequently, if $m$ is prime, $N(\rho)=r\,(q+1-e)$ for each of the six ratios $\rho$ in the $S_3$-orbit of $\rho_1$; this orbit contains $\rho_0=r^4+r^{-4}=(r+1)r^4$.
If $r\equiv1\pmod3$ then $3\mid m$, so for prime $m$ we have $e\in\{0,1\}$, and $e=1$ exactly when $p=3$.
\end{theorem}
\begin{proof}
(i) \emph{Reduction.} Put $u=y/x$. Then $y\in P(L)$ iff $\mathrm{Tr}(xu)=0$. If $x$ is not $\varphi$-fixed, then $\psi=\varphi(x)/x\notin F$ and
$\mathrm{Tr}(x\psi)=\mathrm{Tr}(\varphi x)=0$, so $\{u:\mathrm{Tr}(xu)=0\}=F+F\psi$.
In $K^\times/F^\times$ the exponent $\rho_1$ acts as $u\mapsto\tau(u)\tau^{-1}(u)$, so we need $z=x\,\tau(u)\,\tau^{-1}(u)\in F^\times L$. That is, $\mathrm{Tr}(x\,\tau(u)\,\tau^{-1}(u))=0$.
Take $u=1+s\psi$ with $s\in F$; on $F$ we have $\tau s=s^r$ and $\tau^{-1}s=s^{r^2}$. Using $\mathrm{Tr}(x)=0$, the condition reads $As^r+Bs^{r^2}+Cs^{r+r^2}=0$, where
$A=\mathrm{Tr}(x\tau\psi)$, $B=\mathrm{Tr}(x\tau^{-1}\psi)$ and $C=\mathrm{Tr}(x\,\tau\psi\,\tau^{-1}\psi)$ lie in $F$.
Remove $s=0$ (that is, $y=x$) and put $Z=s^{-r}$, with $s=\infty$ corresponding to $Z=0$. This gives the affine $\mathbb F_r$-linear equation
\[A Z^r+BZ+C=0,\qquad Z\in F.\]
(ii) \emph{Two identities.} Write $\xi_j=\tau^j(x)$ for $j\bmod 9$, with $\xi_j+\xi_{j+3}+\xi_{j+6}=0$. Then $A$, $B$, $C$ are rational functions of $\xi_0,\dots,\xi_5$
with integer coefficients and monomial denominators, and $\tau$ acts as the index shift. Let $D$ be the Dickson matrix of $f(Z)=AZ^r+BZ$. Exact computation in Sage gives
\[N_{F/\mathbb F_r}(A)+N_{F/\mathbb F_r}(B)=0\quad\text{and}\quad \mathrm{adj}(D)\,(C,C^r,C^{r^2})^{T}=0.\]
After clearing the (nonzero) denominators, both are identities between polynomials with integer coefficients that hold over $\mathbb Q$. They therefore hold in every characteristic.
They were also re-checked directly over $\mathbb F_2,\mathbb F_3,\mathbb F_5,\mathbb F_7$.
(iii) \emph{Non-degeneracy.} The numerator of $A$ factors as $(\xi_1\xi_3-\xi_0\xi_4)(\xi_1^2+\xi_1\xi_4+\xi_4^2)$, and the denominator is $-\xi_1\xi_4\xi_7\ne0$.
The second factor vanishes iff $\psi$ is a root of $t^2+t+1$. The first vanishes iff $\tau x=cx$ with $c\in F$, which gives $\varphi x=N_{F/\mathbb F_r}(c)\,x$.
In both cases $x$ is $\varphi$-fixed. So $A\ne0$, hence $N(B)=-N(A)\ne0$ and $\mathrm{rank}\,D=2$.
(iv) \emph{Counting.} $N(-B/A)=N(-1)\,N(B)/N(A)=(-1)^{1+r+r^2}\cdot(-1)=1$, so the kernel of $f$ has exactly $r$ elements and $|\mathrm{Im}\,f|=r^2$.
Since $\mathrm{adj}(D)\,D=\det(D)\,I=0$, each row $(w_0,w_1,w_2)$ of $\mathrm{adj}(D)$ annihilates $(b,b^r,b^{r^2})$ for every $b\in\mathrm{Im}\,f$.
A nonzero row gives a nonzero linearised polynomial of degree at most $r^2$. Its zero set therefore has at most $r^2$ elements and contains $\mathrm{Im}\,f$, so it equals $\mathrm{Im}\,f$.
By (ii), $-C\in\mathrm{Im}\,f$, so there are exactly $r$ solutions $Z$, that is, $r$ partners $y$.
(v) \emph{Fixed points.} The $\varphi$-fixed points of $P(L)$ are the eigenlines of the $F$-linear map $\varphi|_L$, whose eigenvalues are cube roots of unity in $F$.
The eigenvalue $1$ needs $x\in F$, and then $\mathrm{Tr}(x)=3x$ vanishes only if $p=3$.
\begin{itemize}
\item If $r\equiv2\pmod3$, then $F$ has no primitive cube root of unity, so $e=0$.
\item If $p=3$, then $\varphi|_L$ is unipotent with fixed line $F\cdot1$, so $e=1$.
\item If $r\equiv1\pmod3$, then $\varphi|_L$ has the two eigenvalues $\omega,\omega^2$ ($\omega\in F$ a primitive cube root of unity), so $e=2$.
\end{itemize}
For a fixed point:
\begin{itemize}
\item If $\varphi x=\omega x$, then $K$ splits into $\varphi$-eigenlines, which $\tau$ preserves, so $\tau x=cx$ and $\tau^{-1}x=c'x$ with $c,c'\in F$, and $\{u:\mathrm{Tr}(xu)=0\}=F+Fx$.
For $u=\beta_0+\beta x$ the condition becomes $3cc'\beta^{r+r^2}x^3=0$, i.e.\ $\beta=0$, i.e.\ $y=x$.
\item If $p=3$, then $L=F+Fv$ with $\varphi v=v+\alpha$, $\alpha\in F^\times$. For $y=\beta_0+\beta v$, the condition reduces to $\mathrm{Tr}(v^2)\,\beta^{r+r^2}=0$ with $\mathrm{Tr}(v^2)=-\alpha^2\ne0$, so again $y=x$.
\end{itemize}
Summing over $x$ gives $N(\rho_1)$ when $m$ is prime, because then every nonzero difference is a unit.
The maps $(a,b,c)\mapsto(b,a,c)$ and $(a,b,c)\mapsto(a,c,b)$ realise $\rho\mapsto1-\rho$ and $\rho\mapsto1/\rho$ and preserve the unit conditions, so $N$ is constant on the $S_3$-orbit.
\end{proof}
The theorem was also checked by independent brute force for $r\le13$; for $q=8,27,512,1331$ it was checked on all six ratios. An independent AI review accepted it after the fixes included here.

\emph{Algebraic form.} The six ratios of the orbit are roots modulo $m$ of $x^3-3x+1$ and $x^3-3x^2+1$. For prime $m$ they are all of the roots;
for composite $m$ these cubics can have more roots, by the Chinese remainder theorem. The first is the minimal polynomial of $2\cos(2\pi/9)=\zeta_9+\zeta_9^{-1}$, with $\zeta_9\mapsto r$.
Its root set is invariant under $x\mapsto1/(1-x)$, the inverse of the order-3 map $s\mapsto1-1/s$ of Theorem~\ref{thm:curve}(a).

\section{The resonance and the cube excess}\label{sec:resblind}
For the six resonant ratios $M(\rho)\approx r\mu$. If all other ratios behaved as in Conjecture~\ref{conj:var}, then $\mathrm{Var}\,M/\mathbb E\,M$ would exceed 1 by about
\[\Delta=6(r-1)^2\mu/n_{\rm good},\qquad n_{\rm good}=\prod_{\ell^e\|m}\ell^{e-1}(\ell-2),\quad \mu=\frac{(q+1)q(q-1)}{3k\,n_{\rm good}}\prod_{\ell\mid m}\frac{(\ell-1)(\ell-2)}{\ell^2}.\]
This is a heuristic, not a theorem; it decays like $q^{-1/3}$.

\emph{Registration.}
\begin{itemize}
\item The value of $\Delta$ was registered before computation only for $q=12167$.
\item The statements "the six largest $M(\rho)$ are roots of the cubics" and "$M/\mu\in[0.8r,1.2r]$" were registered before computation for $q=4913,6859$. For $q=15625$ the registration was "the cubic orbit is among the twelve largest, with $M/\mu\in[20,30]$".
\item All other entries of Table~\ref{tab:res} are retrospective.
\end{itemize}
The plan file and its SHA-256 ledger are released. Registration times are self-attested.
\begin{table}[h]\centering\small
\begin{tabular}{rrrrrl}\toprule $q$&$k$&$\Delta$ formula&observed&top-6 $M/\mu$ ($r$)&status\\ \midrule
125&3&0.089&0.101&4.91 (5)&retrospective\\
343&3&0.143&0.180&7.19 (7)&retrospective\\
512&9&0.021&0.039&8.03 (8)&retrospective\\
729&6&0.032&0.042&8.86 (9)&retrospective\\
1331&3&0.0500&0.0534&11.0 (11)&retrospective (local computation)\\
2197&3&0.0873&0.0877&-- (13)&retrospective\\
4096&12&0.020&0.022&15.93 (16)&retrospective\\
4913&3&0.0368&0.0377&17.01 (17)&top-6, $M/\mu$ registered\\
6859&3&0.0631&0.0635&18.94 (19)&top-6, $M/\mu$ registered; $\Delta$ retrospective\\
12167&3&0.0281&0.0281&22.99 (23)&\textbf{$\Delta$ registered (blind)}; top-6, $M/\mu$ retrospective\\
15625&6&0.0261&0.0261&24.95 (25)&top-6, $M/\mu$ registered; $\Delta$ not registered\\
19683&9&0.0077&0.0092&27.01 (27)&exploratory ($p=3$)\\
\bottomrule\end{tabular}
\caption{The resonance and the cube excess. Wherever the six largest $M(\rho)$ were recorded, they are roots of the two cubics and $M/\mu\approx r$.
Non-cube controls ($q=8192,14641,16384$) give ratio$/3k\in[0.9997,1.0024]$.}\label{tab:res}\end{table}

For $k=3$ and $q\ge1331$, the formula gives $94$--$100\%$ of the observed excess. In the other rows of Table~\ref{tab:res} it gives $55$--$100\%$ (lowest at $q=512$).
Two earlier registered readings failed at $q=6859$: a $q^{-1/2}$ decay band, and a $p\equiv1\pmod3$ effect. Registered tests at $q=24389$ and $q=32768$ are pending.
For $k=6,9,12$ a second, weaker orbit appears, with $M/\mu=2.77,\,3.91,\,4.00,\,2.98$ at $q=729,4096,15625,19683$. Its ratios are not roots of integer
polynomials of degree $\le4$ with coefficients bounded by 4, and they are not of the form $p^a\pm p^b\pm p^c$. We have not identified this orbit.
At $q=19683$ it contributes only about $10\%$ of the residual $0.0015$; most of the residual comes from several hundred ratios with $M/\mu\approx2$.
Weaker levels with $M/\mu\approx2$ also occur for $k=3$ ($q=4913,12167$) and for the non-cube control $q=16384$. They are also unexplained.
